%%
% 绪论
%%

\chapter{绪论}
\label{cha:introduction}

\section{研究背景}
\label{sec:background}

近年来，生成式人工智能已经进入知识服务、科研辅助与业务分析等场景，企业信息系统的交互方式也随之发生变化\cite{feuerriegel2024generative}。不过，大语言模型在事实一致性和部署可靠性方面仍存在明显风险，幻觉问题依然是当前落地过程中的核心障碍\cite{zhang2025hallucination}。我国《生成式人工智能服务管理暂行办法》进一步将训练数据、生成内容与服务安全纳入规制框架，这意味着企业应用必须同时满足能力建设与合规约束\cite{genai_measures2023}。

从国内数字化建设环境看，《“十四五”国家信息化规划》已将行业信息资源建设和数据要素流通列为持续推进的重点\cite{informatization_plan2021}。国家数据局发布的《全国数据资源调查报告（2024年）》也表明，数据生产与开发利用规模仍在扩大\cite{datareport2024}。与公开互联网场景相比，企业私有知识更强调权限边界、版本更新和来源可追溯性，相关知识管理研究也指出，知识服务的有效性取决于知识组织、权限治理与业务流程的协同\cite{oleary2024km}。

检索增强生成（Retrieval-Augmented Generation, RAG）的基本思想，是在生成前先完成受控检索，再将候选证据组织为模型可利用的上下文。向量空间模型为基于文本表示的相似度匹配奠定了理论基础\cite{salton1975vectorspace}。BM25 所代表的概率检索框架则长期构成关键词检索的重要基线\cite{robertson2009bm25}。中文基准研究表明，RAG 系统性能不仅受召回结果影响，也受生成链路与评测设计共同制约\cite{wang2024crudrag}。国内关于检索粒度的综述进一步指出，候选切分和证据重组会直接影响生成稳定性\cite{raggranularity2025}。

然而，从“可验证的 RAG 原型”到“可落地的企业知识系统”之间仍存在明显的工程鸿沟。真实业务环境下的系统不仅要回答问题，还要承担用户认证、组织隔离、文件上传、断点续传、文档解析、异步向量化、权限过滤、流式对话和历史记录管理等职责。只有将上述能力纳入统一的软件系统中，RAG 才能真正成为企业内部知识管理工具，而不是停留在演示层面的模型能力展示。基于这一现实需求，本文面向企业私有文档场景，设计并实现了一套基于 RAG 的智能知识管理系统，并围绕其需求分析、系统设计与关键实现过程展开研究。

\section{研究意义}
\label{sec:significance}

本文的研究意义首先体现为一个面向企业私有文档场景的软件系统工件。围绕文档来源分散、权限结构复杂和知识更新频繁等现实约束，系统构建了从文件接入、文本解析、向量索引、混合检索到对话生成的完整处理链路，在保留私有文档边界的前提下提供面向内部知识的检索与问答能力。这一工件对应的并不是抽象算法演示，而是企业知识服务最基础的接入、组织与查询需求。

本文也不将 RAG 处理为孤立算法模块，而是将用户管理、组织标签、多租户隔离、分片上传、异步消息、检索重排序和 WebSocket 流式交互纳入同一系统中统一组织，使检索增强生成从局部处理流程转化为具备部署与维护条件的工程方案。这种组织方式直接影响系统的可部署性、可维护性与可扩展性。

从应用适配角度看，企业知识管理是生成式人工智能的重要落地方向。系统以自然语言交互降低内部文档知识的访问门槛，并通过权限过滤、来源引用和检索降级增强回答过程的可控性。相关设计既服务于本项目，也可为制度检索、内部智能客服和企业知识问答等相近场景提供实现参照。

\section{国内外研究现状}
\label{sec:related_work}

围绕本文选题，现有研究大体可以归纳为三条彼此衔接的线索：其一是 RAG 作为知识增强机制的技术演进，其二是企业知识场景中的工程化探索，其三是从研究原型走向平台级系统时暴露出的实现缺口。沿着这一脉络展开，有助于更准确地说明本文所要处理的问题并不只是检索方法选择，而是企业知识服务在权限、接入和交互层面的整体工程组织。

\subsection{RAG 技术的发展}

从研究演进看，RAG 相关工作已经从“是否引入外部检索”推进到“如何组织检索、重排和生成链路”。跨领域比较研究表明，不同架构的差异并不只来自检索器本身，还来自上下文组织方式和生成编排方式\cite{rahman2025comparative}。面向中文知识密集任务的研究则强调，问题理解、检索选择与证据组织需要在统一框架下协同设计\cite{knowledgeintensiverag2025}。

在垂直领域应用中，RAG 的价值更多体现为证据可核对，而不是单纯生成流畅。中医药标准知识问答实践表明，规范性文档场景能够从检索约束中获得更稳定的答案依据\cite{tcmstandardrag2025}。这说明企业知识系统的重点不应停留在“能否作答”，而应进一步落实到证据是否可定位、结果是否可复核。

\subsection{企业知识场景中的工程化研究}

经典的组织知识管理研究指出，知识创造不是静态存储过程，而是组织内部持续互动与外化的结果\cite{nonaka1994dynamic}。进入文档密集型业务场景后，国内关于机构规范文档知识组织的研究也表明，内部文档若缺少稳定的组织体系，后续检索与问答很难持续落地\cite{normdocs2025}。因此，企业级知识系统的关键不只是模型能力，而是知识组织、权限边界与业务流程之间能否形成稳定耦合。

\subsection{现有研究的不足}

综合现有研究可以发现，RAG 相关工作已经为知识增强生成提供了较完备的方法基础，但对于企业知识平台建设而言，研究重点与工程重点并未完全重合。部分工作更关注检索与生成算法本身，对软件系统中的权限隔离、文件接入与长链路稳定性讨论不足，难以直接转化为平台级实现方案。

与此同时，许多工作默认知识已经完成结构化接入，而对“文档如何上传、如何切分、如何异步入库、如何与业务权限绑定”等实现问题涉及较少。对于面向企业非结构化文档的知识系统而言，这些问题恰恰决定了系统能否稳定运行。已有研究通常强调检索质量优化，但对最终用户交互体验的关注相对不足。企业用户不仅关心答案是否准确，也关心上传能否恢复、会话能否连续、结果是否可追溯以及长文本回答能否实时返回。

基于上述不足，本文将研究重点放在面向企业私有知识场景的 RAG 系统工程实现上。

\section{研究内容}
\label{sec:research_content}

围绕系统的真实实现，本文将研究重点放在企业私有知识服务的完整软件链路上。首先，论文从业务使用方式出发梳理系统边界，明确普通用户与管理员两类角色下的操作差异，将私人空间、组织共享空间与公开文档纳入统一的访问控制框架，为后续设计提供约束基础。

在此基础上，本文实现了文档接入与知识入库链路，使系统能够完成分片上传、断点续传、MinIO 文件合并、Kafka 异步投递、Tika 文档解析、HanLP 细粒度切分以及向量化后写入 Elasticsearch 索引。该部分解决的是企业文档从原始文件转化为可检索知识片段的工程问题，也是后续检索与问答成立的前提。

随后，本文实现了带权限约束的混合检索与会话交互模块。一方面，系统将查询向量召回、关键词匹配、组织标签过滤和二阶段重排序整合到同一检索流程中，并为向量生成失败或检索异常保留纯文本降级路径；另一方面，系统基于 WebSocket 实现流式问答，结合 Redis 中的会话历史与检索结果构造提示词，维护引用编号与文件标识之间的映射关系，使回答过程能够与真实文档来源保持关联。

\section{论文结构安排}
\label{sec:arrangement}

全文共分为六章，章节之间按照“问题提出、需求界定、方案设计、系统实现、运行验证、总结展望”的逻辑展开。

第一章为绪论，介绍课题背景、研究意义、国内外研究现状以及本文的研究内容，并说明论文选题所对应的工程问题。

第二章为需求分析，围绕系统建设目标、业务场景、功能需求和非功能需求划定项目边界，为后续设计与实现提供约束条件。

第三章为系统设计，说明总体架构、核心模块划分、权限模型和数据结构设计。

第四章为系统实现，结合关键代码展开权限隔离、文档处理、混合检索与对话交互等模块的实现分析。

第五章为系统测试，围绕第二章提出的功能需求与非功能需求展开验证，对身份建立、文档接入、检索权限、文件服务和流式交互等关键能力进行系统测试，并据此分析当前版本的成立边界。

第六章为总结与展望，就本文的研究工作进行总结，并对未来的研究方向提出展望。
